\begin{abstract}
We prove a uniform upper bound for inverse-first-power integral means
of normalized univalent disk maps with exponent strictly below \(1/4\).
Consequently, the bounded universal integral-means spectrum satisfies
\(B_b(-1)<1/4\), disproving Kraetzer's conjectured spectrum at \(p=-1\).
\end{abstract}
\section{Introduction}\label{sec:introduction}

For a univalent function, the derivative can become large or small near
the boundary even when the image domain is bounded. Integral means
measure the rate of this concentration. Their universal spectrum records
the largest possible growth exponent over a class of conformal maps;
it brings questions about derivatives, coefficients, and the geometry
of image boundaries into a common framework.

Let \(\DD=\{z\in\C:|z|<1\}\), and let \(\SSS\) consist of all
injective holomorphic maps \(f:\DD\to\C\) with \(f(0)=0\) and
\(f'(0)=1\). Write \(\SSS_b\) for the subclass of maps with bounded
image. The bound on the image may depend on \(f\), and no boundary
regularity is imposed. For \(p\in\R\) and \(0\le r<1\), define
\begin{align*}
 M_p[f'](r)&=\frac1{2\pi}\int_{-\pi}^{\pi}|f'(re^{i\theta})|^p\,d\theta,\\
 \beta_f(p)&=\limsup_{r\uparrow1}
 \frac{\log M_p[f'](r)}{\log(1/(1-r))},
 &B_b(p)&=\sup_{f\in\SSS_b}\beta_f(p).
\end{align*}
The derivative of a univalent map never vanishes, so all real powers
are defined. The supremum in \(B_b\) is taken after each map's limiting
exponent. Multiplying a map by a nonzero constant leaves that exponent
unchanged, so normalization of the derivative does not change the
bounded-class spectrum.

\subsection{The conjecture and the uniform bound}

Kraetzer's experiments with conformal maps onto bounded domains led
to the prediction
\begin{equation}\label{eq:kraetzer}
 B_b(p)=
 \begin{cases}
 p^2/4,&|p|\le2,\\
 |p|-1,&|p|\ge2.
 \end{cases}
\end{equation}
The quadratic part is stated in his 1996 paper and explicitly
restated in his thesis \cite{Kraetzer,KraetzerThesis}; the full
quadratic--linear formulation also appears in Beliaev's discussion
of numerical spectra \cite{Beliaev}. We prove a strict upper bound
at the negative interior exponent \(p=-1\).

\begin{theorem}[Uniform inverse-first-power bound]\label{thm:main}
There exist constants \(0<\varepsilon<1/4\) and \(C<\infty\) such that
\begin{equation}\label{eq:uniform}
 M_{-1}[f'](r)\le C(1-r)^{-1/4+\varepsilon}
 \qquad\text{for every }f\in\SSS\text{ and every }1/2\le r<1.
\end{equation}
Consequently \(B_b(-1)<1/4\), and Kraetzer's conjectured identity
\eqref{eq:kraetzer} is false.
\end{theorem}

Both constants are independent of the map, including its image size.
The proof gives existence of a positive gap but does not evaluate it
numerically. The conclusion concerns the inverse first power; it does
